\chapter{Governance}
\label{ch:governance}

This chapter explains how \texttt{gert.ops} wires into gert's core governance layer and extends
it with role-based gating, Kit-specific policy primitives, and audit trail generation. For the
core governance model (allowlists, denylists, env-var blocking, output redaction), see the
\emph{gert v2 Design Document}, §11. This chapter covers the \texttt{gert.ops}-specific
extensions to that model.

% ─────────────────────────────────────────────────────────────────────────────
\section{Core Governance Hooks}
\label{sec:governance-core-hooks}
% ─────────────────────────────────────────────────────────────────────────────

The gert core governance layer evaluates policy at two points in the execution pipeline:

\begin{description}
  \item[Pre-fork checkpoint] Runs immediately before any subprocess is started. Checks the
    command allowlist/denylist and strips blocked environment variables.
  \item[Post-capture checkpoint] Runs after a step completes and output is captured. Applies
    output redaction rules before writing to trace or run variables.
\end{description}

\texttt{gert.ops} adds a third checkpoint:

\begin{description}
  \item[Pre-step role checkpoint] Runs before any step executes. For steps with
    \texttt{requires\_role}, verifies that the current executor holds the required role.
    If the role check fails, the step is blocked and a \texttt{governance/roleBlocked} trace
    event is written.
\end{description}

% ─────────────────────────────────────────────────────────────────────────────
\section{Role-Based Gating}
\label{sec:governance-role-gating}
% ─────────────────────────────────────────────────────────────────────────────

\texttt{gert.ops} steps support three role-gating modifiers. These are evaluated at the
pre-step role checkpoint.

\subsection{\texttt{requires\_role}}

The executing user must hold exactly the specified role.

\begin{minted}{yaml}
- step:
    id: approve_rollback
    type: ops.manual
    title: "Authorize rollback"
    requires_role: change-manager
    instructions: |
      Review the rollback plan and confirm authorization.
\end{minted}

\subsection{\texttt{requires\_any\_role}}

The executing user must hold at least one of the listed roles.

\begin{minted}{yaml}
- step:
    id: triage_check
    type: ops.manual
    title: "Initial triage"
    requires_any_role:
      - dri
      - incident-commander
    instructions: |
      Perform initial triage and document findings.
\end{minted}

\subsection{\texttt{requires\_all\_roles}}

The executing user must hold all of the listed roles simultaneously. Rare in practice; used
when a single person must be both DRI and change-manager (e.g., emergency change procedures).

\begin{minted}{yaml}
- step:
    id: emergency_authorize
    type: ops.manual
    title: "Emergency change authorization"
    requires_all_roles:
      - dri
      - change-manager
    instructions: |
      As both DRI and change manager, authorize this emergency change.
\end{minted}

\paragraph{Role inheritance.} The \texttt{incident-commander} role inherits all permissions of
\texttt{dri}. During an \texttt{ops.incident} wrapper execution, the declared commander may
execute any step gated on \texttt{requires\_role: dri}.

% ─────────────────────────────────────────────────────────────────────────────
\section{Allowlist Enforcement in \texttt{ops.cli} Steps}
\label{sec:governance-allowlist}
% ─────────────────────────────────────────────────────────────────────────────

Command allowlists in \texttt{gert.ops} follow the core model but are evaluated hierarchically:

\begin{enumerate}
  \item \textbf{Runbook-level:} \texttt{meta.governance.allowed\_commands} — applies to all
    steps unless overridden.
  \item \textbf{Step-level:} \texttt{ops.cli.allowed\_commands} — intersected with the
    runbook-level list. A command not in the runbook-level list cannot be unblocked at the
    step level.
  \item \textbf{Denylist:} \texttt{meta.governance.denied\_commands} — always applied first,
    regardless of allowlist.
\end{enumerate}

\begin{minted}{yaml}
meta:
  governance:
    allowed_commands: [kubectl, helm, curl, jq, aws]
    denied_commands: [rm, dd, shutdown, reboot, curl]  # curl denied globally

flow:
  - step:
      id: check_pods
      type: ops.cli
      title: "List running pods"
      command: kubectl
      args: [get, pods, -n, production]
      # Step-level allowlist narrows to kubectl only; curl remains denied.
      allowed_commands: [kubectl]
\end{minted}

When a command is blocked, gert writes a \texttt{governance/commandBlocked} trace event and
fails the step with exit code 126.

% ─────────────────────────────────────────────────────────────────────────────
\section{Environment Variable Blocking}
\label{sec:governance-env-blocking}
% ─────────────────────────────────────────────────────────────────────────────

The \texttt{deny\_env\_vars} list accepts glob patterns. Before any subprocess fork, the
full process environment is scanned and matching variable names are stripped.

\begin{minted}{yaml}
meta:
  governance:
    deny_env_vars:
      - "*_SECRET*"         # AWS_SECRET_ACCESS_KEY, DB_SECRET_PASSWORD, etc.
      - "*_TOKEN"           # GITHUB_TOKEN, VAULT_TOKEN, etc.
      - "AWS_*"             # All AWS credential vars
      - "KUBECONFIG"        # Prevent implicit cluster targeting
      - "DOCKER_HOST"       # Prevent Docker daemon redirection
\end{minted}

\paragraph{Trace event.} When variables are stripped, a \texttt{governance/envVarBlocked}
event is written. The event records the \emph{names} of blocked variables, never the values:

\begin{minted}{yaml}
# Trace event: governance/envVarBlocked
kind: governance/envVarBlocked
payload:
  step_id: "deploy_image"
  var_names: ["AWS_SECRET_ACCESS_KEY", "GITHUB_TOKEN"]
  reason: "matched deny_env_vars pattern"
\end{minted}

% ─────────────────────────────────────────────────────────────────────────────
\section{Output Redaction Patterns}
\label{sec:governance-redaction}
% ─────────────────────────────────────────────────────────────────────────────

Output redaction applies to the captured stdout and stderr of \texttt{ops.cli} steps before
the output is stored in run variables or written to evidence records.

\begin{minted}{yaml}
meta:
  governance:
    redact:
      # Bearer tokens in HTTP responses
      - pattern: "Bearer [A-Za-z0-9._\\-]+"
        replace: "Bearer [REDACTED]"
      # Password in connection strings
      - pattern: ":[^@:]+@"
        replace: ":[REDACTED]@"
      # AWS access keys (20 uppercase alphanum chars)
      - pattern: "(?:AKIA|ASIA)[A-Z0-9]{16}"
        replace: "[AWS-KEY-REDACTED]"
      # kubectl --token flag
      - pattern: "(--token=)[^ ]+"
        replace: "$1[REDACTED]"
\end{minted}

Patterns use RE2 syntax. The \texttt{replace} field may reference capture groups with
\texttt{\$1}, \texttt{\$2}, etc. Rules are applied in declaration order. After all rules
are applied, the redacted output is stored; the raw output is discarded and never written
to the trace.

% ─────────────────────────────────────────────────────────────────────────────
\section{The Audit Trail}
\label{sec:governance-audit-trail}
% ─────────────────────────────────────────────────────────────────────────────

The \texttt{gert.ops} audit trail is the sequence of governance events in the run trace.
The following events are written by the governance layer:

\begin{description}
  \item[\texttt{governance/roleBlocked}] A user attempted to execute a step requiring a role
    they do not hold. Fields: \texttt{step\_id}, \texttt{user}, \texttt{required\_role},
    \texttt{actual\_roles}.
  \item[\texttt{governance/commandBlocked}] A command was blocked by the allowlist or
    denylist. Fields: \texttt{step\_id}, \texttt{command}, \texttt{reason}.
  \item[\texttt{governance/envVarBlocked}] Environment variables were stripped before
    subprocess fork. Fields: \texttt{step\_id}, \texttt{var\_names}.
  \item[\texttt{governance/approvalGranted}] An approval gate was signed. Fields:
    \texttt{step\_id}, \texttt{approved\_by}, \texttt{role}, \texttt{timestamp}.
  \item[\texttt{governance/approvalRejected}] An approval gate was rejected. Fields:
    \texttt{step\_id}, \texttt{rejected\_by}, \texttt{role}, \texttt{reason}.
  \item[\texttt{governance/approvalTimeout}] An approval gate timed out. Fields:
    \texttt{step\_id}, \texttt{action} (\texttt{escalate} or \texttt{abort}).
  \item[\texttt{evidence/submitted}] An evidence record was submitted. Fields:
    \texttt{step\_id}, \texttt{evidence\_label}, \texttt{evidence\_type},
    \texttt{submitted\_by}, \texttt{role}, \texttt{sha256}.
\end{description}

These events, combined with the core \texttt{step/started} and \texttt{step/completed}
events, form a tamper-evident audit log. The audit trail can be projected from the trace
using \texttt{gert kit audit-report <run-id>}.

% ─────────────────────────────────────────────────────────────────────────────
\section{Complete Governance Configuration Example}
\label{sec:governance-full-example}
% ─────────────────────────────────────────────────────────────────────────────

\begin{minted}{yaml}
apiVersion: runbook/v2
kit: ops/v1

meta:
  name: infrastructure-change
  kind: change
  roles:
    dri: "@ops-lead"
    approver:
      - "@infra-change-board"
    change-manager: "@infra-change-board"
  sla:
    approval-timeout: 8h
    escalation-target: "@vp-ops"
  governance:
    allowed_commands:
      - terraform
      - aws
      - kubectl
      - jq
      - curl
    denied_commands:
      - rm
      - dd
      - shutdown
      - reboot
    deny_env_vars:
      - "*_SECRET*"
      - "*_TOKEN"
      - "AWS_SECRET_ACCESS_KEY"
      - "AWS_SESSION_TOKEN"
      - "TF_VAR_*"     # Block all Terraform variable overrides
    redact:
      - pattern: "(password\\s*=\\s*)[^\\s]+"
        replace: "$1[REDACTED]"
      - pattern: "(?:AKIA|ASIA)[A-Z0-9]{16}"
        replace: "[AWS-KEY-REDACTED]"
      - pattern: "Bearer [A-Za-z0-9._\\-]+"
        replace: "Bearer [REDACTED]"

flow:
  - step:
      id: infra_approval
      type: ops.approval
      title: "Infrastructure change approval"
      description: |
        Terraform plan for {{ .environment }} environment.
        Risk level: {{ .risk_level }}
        Ticket: {{ .change_ticket }}
      approvers: ["@infra-change-board"]
      requires_any: true
      timeout: 8h
      on_timeout: escalate
      evidence:
        - type: ops.evidence.attestation
          label: "Infra change board approval"
          required: true
          prompt: "I approve this infrastructure change for {{ .environment }}."

  - step:
      id: apply_terraform
      type: ops.change-request
      title: "Apply Terraform plan"
      change-id: "{{ .change_ticket }}"
      risk: "{{ .risk_level }}"
      rollback:
        runbook: ./rollback-terraform.runbook.yaml
        inputs:
          environment: "{{ .environment }}"
      steps:
        - step:
            id: tf_apply
            type: ops.cli
            title: "Terraform apply"
            command: terraform
            args: [apply, -auto-approve, "{{ .plan_file }}"]
            allowed_commands: [terraform]
            evidence:
              auto: true
              label: "terraform apply output"

  - step:
      id: dri_sign_off
      type: ops.manual
      title: "DRI post-change sign-off"
      requires_role: dri
      instructions: |
        Verify all Terraform resources are in the expected state.
        Check CloudWatch/monitoring for anomalies post-apply.
      evidence:
        - type: ops.evidence.attestation
          label: "DRI post-change confirmation"
          required: true
          prompt: "I confirm the infrastructure change was applied successfully with no anomalies."
\end{minted}
